\documentclass[11pt]{article}
\usepackage{amsmath,amssymb,amsthm}
\newtheorem{lemma}{Lemma}
\newtheorem{theorem}{Theorem}
\title{Interleaved Ranking Votes Need Many M\"obius Exceptions}
\author{}
\date{November 24, 2026}
\begin{document}
\maketitle

\paragraph{Scope and status.}
This is a complete written obstruction to one proposed orientation
repair, not a lower bound for arbitrary admissible heads. The proposal
transports each forward Cauchy denominator separately by a fractional
linear transformation, adds the associated affine numerator correction,
and removes a few exceptional row or column vertices to obtain the
whole-cube reverse denominator budget. For interleaved geometric
ranking values, almost half the vertices must be removed from a single
head. Joint changes to the denominator decomposition, interactions
between tails, and arbitrary local compilers are not covered. Nothing
here resolves either global asymptotic target. The results are not Lean
formalized. The companion script makes exact rational checks of the
identities and finite instances; the universal quantifiers are justified
by the proof below.

\begin{lemma}[The exact grid budget]
Let $P=(p_i)_{i\in I}$ and $D=(d_j)_{j\in J}$ be nonempty finite
lists with $\max P<\min D$, and put $r=|I|$, $s=|J|$, with $r+s>2$.
There is an affine denominator on the Boolean $(r+s)$-cube with
strictly negative slopes, positive on the whole cube, and values
$d_j-p_i$ at $(e_i,e_j)$ if and only if
\[
 \sum_i(p_i-\min P)+\sum_j(\max D-d_j)<\max D-\min P.
 \tag{1}
\]
The same criterion holds if an arbitrary positive common scale is
permitted on the grid values.
\end{lemma}
\begin{proof}
Write the desired denominator as
$\beta-\sum_i u_i z_i-\sum_j v_j w_j$, with all $u_i,v_j>0$.
Its grid values force
$u_i=p_i-p_0$, $v_j=d_0-d_j$, and $\beta=d_0-p_0$ for constants
$p_0<\min P$ and $d_0>\max D$. Its all-one value is
\[
 (\max D-\min P)-\sum_i(p_i-\min P)-\sum_j(\max D-d_j)
 -(r-1)(\min P-p_0)-(s-1)(d_0-\max D).
\]
Positivity therefore implies (1). Conversely, if (1) has slack $\delta$,
set $p_0=\min P-\eta$ and $d_0=\max D+\eta$ for
$0<\eta<\delta/(r+s-2)$. All slopes are then strictly negative and
the all-one value is positive. It is the minimum on the Boolean cube.
Positive scaling preserves the criterion.
\end{proof}

\begin{lemma}[A cross-ratio contribution]
Suppose $0\leq x,y<1$, $x+y<1$, and
\[
 \frac{xy}{(1-x)(1-y)}=\rho.
\]
Then $x+y\geq\rho$. If $R>1$, $U,V>0$ and
$0<u,v\leq\min(U,V)/R$, then
\[
 \frac{(U-u)(V-v)}{(U+v)(V+u)}
 \geq\left(\frac{R-1}{R+1}\right)^2.
 \tag{2}
\]
The same bound holds whenever $U,V\geq R\max(u,v)>0$.
\end{lemma}
\begin{proof}
Since $(1-x)(1-y)-xy=1-x-y$,
$1-x-y=(1-\rho)(1-x)(1-y)\leq1-\rho$.
For the second assertion each factor in the numerator is at least
$1-1/R$ times its large variable, and each denominator factor is at
most $1+1/R$ times that variable. The final hypotheses give the same
four inequalities.
\end{proof}

\begin{theorem}[Linear exception cost for interleaved orders]
Fix integers $k\geq3$ and a real $R\geq6$. Consider the forward
denominator table
\[
 B_{ij}=u_i+v_j,\qquad
 u_i=R^{2i},\quad v_j=R^{2j+1},\qquad 0\leq i,j<k.
\]
Retain row indices $I$ and column indices $J$, deleting at most $k-3$
vertices in total. Let $\phi$ be any nonconstant real fractional linear
map defined at all $-u_i$ and $v_j$ for retained indices. If
\[
 \phi(v_j)-\phi(-u_i)>0\qquad(i\in I,j\in J),
\]
then these transformed grid denominators cannot be extended to a
whole-cube positive affine denominator with weakly nonpositive slopes.
In particular they cannot satisfy strict reverse admissibility.
Consequently this headwise M\"obius repair must delete at least $k-2$
vertices from this table before strict admissibility can be possible.
\end{theorem}
\begin{proof}
The deletion hypothesis says $|I|+|J|\geq k+3$, and hence both retained
sets have at least three members.

First suppose $\phi$ is affine or its pole is outside the interval
from $-\max_{i\in I}u_i$ to $\max_{j\in J}v_j$. Positivity of every
transformed difference forces $\phi$ to be increasing on this interval.
Put $U=\max_{i\in I}u_i$, $V=\max_{j\in J}v_j$ and
$M=\min(U,V)$. There are at least two retained row values and at least
two retained column values at most $M/R$.
To check this explicitly, put $a=\max I$, $b=\max J$. If $a\leq b$,
then $M=R^{2a}$; all retained rows except $a$ qualify, and the number
of retained columns $j\leq a-1$ is at least
$|J|+a-k\geq|I|+|J|-k-1\geq2$.
If $b<a$, then $M=R^{2b+1}$; all retained columns except $b$ qualify,
and there are at least $|I|+b+1-k\geq|I|+|J|-k\geq3$ qualifying
rows.

Normalize the transformed endpoints $\phi(-U),\phi(V)$ to $0,1$ by
an increasing affine map. For each qualifying row and column pair let
$x$ be its normalized row value and let $y$ be one minus its normalized
column value. Fractional linear invariance of the cross-ratio gives
\[
 \frac{xy}{(1-x)(1-y)}
 =\frac{(U-u_i)(V-v_j)}{(U+v_j)(V+u_i)}.
\]
The preceding lemma yields
$x+y\geq((R-1)/(R+1))^2$. Choose two distinct qualifying rows and
two distinct qualifying columns and pair them arbitrarily. Their four
contributions alone make the normalized left side of (1) at least
\[
 2\left(\frac{R-1}{R+1}\right)^2
 \geq\frac{50}{49}>1.
 \tag{3}
\]
Thus the budget fails even with weak slopes.

It remains to consider a pole between the retained negative and
positive blocks. Positivity forces $\phi$ to be decreasing on each
block. Put $u=\min_{i\in I}u_i$, $v=\min_{j\in J}v_j$ and
$M=\max(u,v)$. There are at least two retained rows and two retained
columns whose values are at least $RM$. Indeed, writing
$a=\min I$, $b=\min J$, if $a\leq b$ then $M=R^{2b+1}$.
All retained columns except $b$ qualify, and there are at least
$|I|-b-1\geq|I|+|J|-k-1\geq2$ qualifying rows. If $b<a$, then
$M=R^{2a}$; all rows except $a$ qualify, and at least
$|J|-a\geq|I|+|J|-k\geq3$ columns qualify.

Normalize $\phi(-u),\phi(v)$ to $0,1$. These are now respectively
the smallest transformed row value and largest transformed column
value. For a qualifying pair with original values $U=u_i,V=v_j$,
the same cross-ratio formula is
\[
 \frac{xy}{(1-x)(1-y)}
 =\frac{(U-u)(V-v)}{(U+v)(V+u)}.
\]
The final part of the lemma and two disjoint pairs again give (3).

Finally no other pole location is possible. For
$\phi(t)=(\alpha t+\beta)/(\gamma t+\delta)$,
\[
 \phi(d)-\phi(p)=
 \frac{(\alpha\delta-\beta\gamma)(d-p)}
 {(\gamma d+\delta)(\gamma p+\delta)}.
\]
A pole inside either retained block changes the sign of one denominator
factor while the other block has fixed sign. The transformed grid
differences could not then all be positive. This exhausts the cases.
\end{proof}

\paragraph{Interpretation.}
The ranks in this example alternate between row and column vertices.
Such an order already occurs in a paired-order representation of a
bilevel graph consisting of disjoint directed edges. Affine row and
column corrections can repair the transported numerators, but do not
change the failed denominator budget. Therefore a theorem permitting
only a constant number of exceptional vertices per transported vote
cannot hold for all rankings. The result does not force a comparable
exception count after replacing several votes jointly by a different
family of heads.
\end{document}
